\section{Introduction}

This document describes the \textbf{ntiny} core one module at a time, bottom-up:
the memory leaves first, then the caches, the bus, the load/store and atomic
units, the core control logic, and finally how they compose into \sig{core\_top}
and the SoC. Each module gets a block diagram, a port summary, a plain-language
walk-through, and---where it matters---a note on how it behaves when memory does
\emph{not} answer in a single cycle.

\subsection{The pipeline, in one breath}

ntiny is a classic in-order pipeline. Instructions flow through five stages:

\begin{center}
\sig{IF} (fetch) $\rightarrow$ \sig{ID} (decode/align) $\rightarrow$
\sig{IE} (execute) $\rightarrow$ \sig{IMEM} (memory) $\rightarrow$
\sig{IWB} (writeback).
\end{center}

A fetch front-end (PC logic, instruction cache, a fetch buffer, and a compressed
aligner) feeds a decoder; the execute stage resolves branches and runs the ALU;
the memory stage issues data accesses through a shared data bus; writeback
commits results and CSR side effects. Forwarding paths and a central
\sig{hazard\_unit} keep the pipeline correct without software-visible stalls.

\subsection{The one idea that explains the hard bugs}

\latencynote{The original pipeline was \emph{statically scheduled} around a
memory that always answers in exactly one cycle. With single-cycle memory, ``the
load's data is on the bus exactly when the consumer needs it,'' so no dynamic
interlock is required---correctness falls out of the fixed schedule. The moment
memory latency becomes \emph{variable} (a real cache miss, \sig{+RAM\_RANDOM\_DELAY},
a future AXI/DRAM slave), that invariant breaks everywhere a memory result is
produced or consumed: forwarding, branch/JALR resolution, the writeback commit
gate, the fetch in-flight tracking, the bus response routing, and the cache
itself. Each such site needs an explicit ``data-not-ready-yet'' handshake. This
is why ``just more wait time'' is deceptively invasive---it is not one stall
point, it is the removal of a global timing assumption. Throughout this document,
\textbf{Latency notes} flag exactly where that assumption lived and what replaced
it.}

\subsection{How to read the diagrams}

Every diagram uses one consistent visual vocabulary so shapes always mean the
same thing:

\begin{center}
\begin{tabular}{>{\raggedright}p{3.2cm} l}
\toprule
\textbf{shape / colour} & \textbf{meaning} \\
\midrule
yellow rectangle      & combinational logic (decode, mux net, ALU op) \\
blue rectangle        & register / flip-flop (pipeline reg, latch, counter) \\
green stadium         & a state machine (state register + transitions) \\
grey hexagon          & multiplexer / select \\
purple cylinder       & memory array (SRAM, cache tag/data arrays) \\
white box             & an instantiated sub-module \\
pale boxes (in/out)   & module ports, grouped left (inputs) / right (outputs) \\
\midrule
solid arrow           & data / datapath signal \\
dotted arrow          & control / handshake / stall / freeze \\
\bottomrule
\end{tabular}
\end{center}

Diagrams are deliberately \emph{sparse}: node labels are short signal or block
names, and all the detail lives in the prose of each section. The diagram sources
are Mermaid (\texttt{docs/microarch/diagrams/*.mmd}); rebuild the figures and this
PDF with \texttt{python3 docs/microarch/build.py}.
